% TOP_PROBLEM: {"id": 3082, "title": "Inequality of the means", "category_id": 6, "category": {"id": 6, "name": "geometry", "display_name": "Geometry", "description": "Euclidean and non-Euclidean geometry, geometric structures.", "slug": "geometry", "order_index": 6, "created_at": "2026-07-31T15:26:25.671Z"}, "status": "open", "problem_number": "OPG-36901", "set_id": 12, "statusQualification": "Ordinary congruent packing allows rotations; the stronger axis-parallel variant is identified rather than silently imposed."}
% FIRST_POSTED: 2026-10-01
% ENTRY_KIND: statement_only
% UnsolvedMath ID 3082; source catalogue code OPG-36901
% !TeX program = lualatex
% Definition and sources only; this file is not an LLM solution attempt.
\documentclass[11pt,a4paper]{article}
\usepackage[margin=25mm]{geometry}
\usepackage{fontspec}
\setmainfont{lmroman10-regular.otf}[BoldFont=lmroman10-bold.otf,
 ItalicFont=lmroman10-italic.otf,BoldItalicFont=lmroman10-bolditalic.otf]
\usepackage{amsmath,amssymb,mathtools}
\usepackage{unicode-math}
\setmathfont{latinmodern-math.otf}
\AtBeginDocument{\let\setminus\smallsetminus}
\usepackage{xurl,hyperref}
\hypersetup{colorlinks=true,urlcolor=blue}
\setcounter{secnumdepth}{0}
\setlength{\parindent}{0pt}
\setlength{\parskip}{5pt}
\setlength{\emergencystretch}{3em}
\begin{document}
\section*{Inequality of the means}\label{3082}
{\small UnsolvedMath ID 3082 · OPG-36901 · Geometry · OpenGarden}
\subsection{Definitions and mathematical statement}
Let $n\ge1$ and let $a_1,\ldots,a_n$ be arbitrary positive real numbers.
Set $s=\sum_{i=1}^n a_i$, let
$B=\prod_{i=1}^n[0,a_i]$ be a rectangular box, and let $Q=[0,s]^n$.
A packing of $m$ congruent copies of $B$ in $Q$ is a family of Euclidean
isometries $T_1,\ldots,T_m$ such that $T_j(B)\subseteq Q$ for every $j$
and the interiors of different copies are disjoint. Boundary contacts
are allowed; boxes may be moved and rotated but not cut or deformed.

Does such a packing always exist with $m=n^n$? The quantifiers range
over every dimension and every choice of side lengths. Permuting the
side lengths changes only the presentation of the same box, not the
available shape. The definition uses the ordinary unrestricted meaning
of geometric packing in the catalogue question. Requiring every edge to
be parallel to a coordinate axis would be a stronger version; the source's
illustrated constructions use such orientations.

A necessary volume comparison is
\[
 n^n\prod_{i=1}^n a_i\le\left(\sum_{i=1}^n a_i\right)^n,
\]
which is exactly the arithmetic--geometric mean inequality after taking
$n$th roots. The problem asks for actual disjoint placement, not just this
numerical inequality. When all lengths equal $a$, the $n\times\cdots\times n$
grid tiles the target cube. In the unequal case, unused space is permitted
and the copies need not cover the entire cube.
\subsection{Short English statement}
Can $n^n$ congruent boxes with side lengths $a_1,\ldots,a_n$ always fit into a cube of side $a_1+\cdots+a_n$?
\subsection{Sources}
\begin{itemize}
\item[S1] ULAM.AI and the UnsolvedMath Contributors, supplied problems.json, record 3082 (OPG-36901), OpenGarden collection. Catalogue identity and source transcription; CC BY 4.0.
\url{https://huggingface.co/datasets/ulamai/UnsolvedMath}.
\item[S2] Open Problem Garden, Inequality of the means. Original problem and discussion; the mathematical formulation below is expanded from the supplied source record.
\url{http://www.openproblemgarden.org/op/inequality_of_the_means}.
\item[S3] D. Bar-Natan, Inequality of the Means, author’s illustrated formulation and constructions.
\url{https://www.math.toronto.edu/~drorbn/projects/ArithGeom/}.
\end{itemize}
\end{document}
